% ============================================================================
%  WAVE 4, TARGET 1c: EXACT MINIMAX SAMPLE COMPLEXITY FOR THE 3-WORLD PROBLEM
%  Closes the residual κ/4 gap from wave3_tight_and_computability.tex (T1b).
%
%  Resolution: κ(α) is EXACTLY optimal among all (possibly randomized) rules;
%  the floor 4 n_opt is a valid but LOOSE pairwise bound; κ/4 > 1 at every
%  finite α is an IMPOSSIBILITY theorem (certification at 4 n_opt cannot hold).
%
%  Validator: theory_v2/val_tight_constants.py  (section T1c)
%
%  REVISION NOTE (proof correction): Step 1 previously invoked MLR / Karlin-Rubin to
%  claim the symmetric two-sided set is the P_0-constrained maximizer of P_{+Delta}
%  mass. That is wrong -- the likelihood ratio is monotone increasing, so Neyman-Pearson
%  gives the ONE-SIDED upper tail. The constant kappa(alpha) is unaffected; the argument
%  is now the correct budget-split Neyman-Pearson one (p + q <= alpha, minimized at
%  p = q = alpha/2). The file also used z_alpha and z_{1-alpha/2} under two conflicting
%  conventions; it now uses z_{1-x} := Phi^{-1}(1-x) throughout.
% ============================================================================
\providecommand{\TV}{\mathrm{TV}}
\section{Exact minimax sample complexity for honest frontier abstention}\label{sec:t1c-closure}

Wave~3 (Theorem~\ref{thm:minimax-opt}, Prop.~\ref{prop:t1b-3world}--\ref{prop:t1b-achieve})
established: (i)~$4n_{\mathrm{opt}}$ is a valid \emph{all-rules} lower bound once a third
frontier world $\mu=0$ must abstain at level $\alpha$; (ii)~the symmetric band certifies at
$n=\kappa(\alpha)\,n_{\mathrm{opt}}$ with
$\kappa(\alpha)=\bigl(1+z_{1-\alpha/2}/z_{1-\alpha}\bigr)^2$, which lies in $[4,\approx5.2]$
\emph{for $\alpha\le0.10$} and decreases to $4$ as $\alpha\downarrow0$ --- the range is not
universal: $\kappa(0.20)=6.36$, and $\kappa$ is increasing in $\alpha$ on $(0,\tfrac12)$;
(iii)~the residual $\kappa/4\in[1,1.30]$ was left open for arbitrary randomized rules.
This section closes (iii).

\subsection{The three-world certifying problem}

Fix $\sigma>0$, $\Delta>0$, $\alpha\in(0,\tfrac14)$, and write $v=\sigma^2/n$,
$z_{1-x}:=\Phi^{-1}(1-x)$ for $x\in(0,1)$ --- so $z_{1-\alpha}=\Phi^{-1}(1-\alpha)$ and
$z_{1-\alpha/2}=\Phi^{-1}(1-\alpha/2)$, both positive, and the earlier shorthand $z_\alpha$ for
$\Phi^{-1}(1-\alpha)$ is retired because it collided with this subscript convention ---
and $n_{\mathrm{opt}}=(\sigma/\Delta)^2 z_{1-\alpha}^2$.  Observe
$\bar X\sim\mathcal N(\mu,v)$ under world $\mu\in\{-\Delta,0,+\Delta\}$ (sufficient statistic
for the $n$-sample Gaussian location model).  A (possibly randomized) \emph{certifying rule} maps
$\bar X$ to $\{\textsc{adapt},\textsc{freeze},\textsc{abstain}\}$ and is
\emph{$\alpha$-certifying} if
\[
  \mathrm{miss}_+:=\Pr_{+\Delta}(\text{not \textsc{adapt}})\le\alpha,\quad
  \mathrm{miss}_-:=\Pr_{-\Delta}(\text{not \textsc{freeze}})\le\alpha,\quad
  \mathrm{commit}_0:=\Pr_{0}(\text{not \textsc{abstain}})\le\alpha,
\]
with the usual wrong-commit rates subsumed by the miss bounds.  The \emph{minimax certifying
sample complexity} is
\[
  n^\star_{3}(\Delta,\sigma,\alpha)
  \;=\;\min\bigl\{\,n\in\mathbb N:\ \exists\ \text{an $\alpha$-certifying rule on }
  \{\pm\Delta,0\}\bigr\}.
\]

\begin{theorem}[Exact minimax for the three-world problem]\label{thm:t1c-exact}
For every $\sigma>0$, $\Delta>0$, and $\alpha\in(0,\tfrac14)$,
\[
  \boxed{\;
  n^\star_{3}(\Delta,\sigma,\alpha)
  \;=\;\kappa(\alpha)\,n_{\mathrm{opt}}(\Delta,\sigma,\alpha)
  \;=\;\Bigl(\tfrac{\sigma\,(z_{1-\alpha}+z_{1-\alpha/2})}{\Delta}\Bigr)^2\; }.
\]
The symmetric band ``\textsc{adapt} iff $\bar X>\tau$; \textsc{freeze} iff $\bar X<-\tau$;
else \textsc{abstain}'' with $\tau=z_{1-\alpha/2}\,\sigma/\sqrt n$ attains this with equality
(all three constraints bind).  Consequently:
\begin{enumerate}
\item[(a)] \textbf{Optimality of $\kappa$.}  No $\alpha$-certifying rule --- deterministic or
randomized --- exists at any $n<\kappa(\alpha)\,n_{\mathrm{opt}}$.
\item[(b)] \textbf{Impossibility of closing $\kappa/4$.}  For every finite $\alpha$,
$\kappa(\alpha)>4$ strictly, so certification at $n=4n_{\mathrm{opt}}$ is \emph{impossible}
even though $n\ge4n_{\mathrm{opt}}$ is a valid \emph{necessary} condition from the
world-$0$ vs.\ world-$+$ pairwise reduction.  The gap $\kappa/4>1$ is intrinsic, not a proof
artifact.
\item[(c)] \textbf{Randomization does not help.}  Every $\alpha$-certifying rule at
$n=\kappa(\alpha)\,n_{\mathrm{opt}}$ may be taken non-randomized and monotone in $\bar X$.
\end{enumerate}
\end{theorem}

\begin{proof}
Write $t=\sqrt v=\sigma/\sqrt n$.

\emph{Achievability (upper bound).}  Take the symmetric band with half-width
$\tau=z_{1-\alpha/2}\,t$.  Then
$\mathrm{commit}_0=2\Phi(-\tau/t)=2\Phi(-z_{1-\alpha/2})=\alpha$ exactly.  In world $+\Delta$,
$\mathrm{miss}_+=\Phi\bigl((\tau-\Delta)/t\bigr)$; feasibility requires
$(\Delta-\tau)/t\ge z_{1-\alpha}$, i.e.\ $\Delta/t\ge z_{1-\alpha}+z_{1-\alpha/2}$.  The smallest $n$
satisfying this is $n=\kappa(\alpha)\,n_{\mathrm{opt}}$, at which $\mathrm{miss}_+=\alpha$ and,
by symmetry, $\mathrm{miss}_-=\alpha$.  Hence $n^\star_{3}\le\kappa(\alpha)\,n_{\mathrm{opt}}$.

\emph{Lower bound (all rules).}  Let $g$ be any $\alpha$-certifying rule at sample size $n$.
By the standard sufficiency of $\bar X$ for $\mu$ in $\mathcal N(\mu,v)$, we may assume $g$
depends only on $\bar X$ (randomization included: a randomized rule is a mixture of
$\bar X$-measurable rules, and all error constraints are linear in the mixture weights).

Define the \emph{commit region} $C\subseteq\mathbb R$ (the set where $g$ outputs
\textsc{adapt} or \textsc{freeze}, possibly with randomization averaged out).  Then
\[
  \mathrm{commit}_0=\Pr_0(\bar X\in C)\le\alpha,\qquad
  \mathrm{miss}_+=\Pr_{+\Delta}(\bar X\notin A_+)\le\alpha,
\]
where $A_+$ is the \textsc{adapt} region.  Since wrong \textsc{freeze} in world $+\Delta$ is a
subset of $\bar X\notin A_+$, the miss bound is the operative constraint on the high side; by
symmetry the low side is identical.

\emph{Step 1 (necessary abstention width; Neyman--Pearson with a split budget).}
\textbf{Correction.}  An earlier version of this step invoked the monotone-likelihood-ratio /
Karlin--Rubin property to claim that the two-sided set
$(-\infty,-\tau]\cup[\tau,\infty)$ maximizes $\Pr_{+\Delta}(\bar X\in C)$ at fixed $\Pr_0(C)$.
That is \emph{false}: $dP_{+\Delta}/dP_0$ is increasing in $x$, so by Neyman--Pearson the
maximizer is the \emph{one-sided} upper tail $[\tau',\infty)$, and the symmetric set is strictly
sub-optimal for that objective.  The conclusion below is nevertheless correct; here is an argument
that actually delivers it.

Let $A_+$ and $A_-$ be the \textsc{adapt} and \textsc{freeze} regions and write
$p:=\Pr_0(A_+)$, $q:=\Pr_0(A_-)$, so that the abstention constraint reads
$\mathrm{commit}_0=p+q\le\alpha$.  These are two \emph{separate} one-sided testing problems that
share a single type-I budget.  Apply Neyman--Pearson to each side individually: among all sets of
$P_0$-measure $p$, the one maximizing $P_{+\Delta}$-mass is the upper tail
$[\,t\,\Phi^{-1}(1-p),\infty)$, so
$\mathrm{miss}_+\le\alpha$ forces
$\Delta/t\ \ge\ z_{1-\alpha}+\Phi^{-1}(1-p)$, and symmetrically
$\Delta/t\ \ge\ z_{1-\alpha}+\Phi^{-1}(1-q)$.
Any rule must satisfy both, i.e.\
\[
  \Delta/t\ \ge\ z_{1-\alpha}+\max\{\Phi^{-1}(1-p),\,\Phi^{-1}(1-q)\}.
\]
Since $\Phi^{-1}(1-\cdot)$ is decreasing, the right-hand side is minimized over the budget set
$\{p+q\le\alpha,\ p,q\ge0\}$ at $p=q=\alpha/2$, giving
\[
  \Delta/t\;\ge\;z_{1-\alpha}+z_{1-\alpha/2}
  \quad\Longleftrightarrow\quad
  n\;\ge\;\Bigl(\tfrac{\sigma\,(z_{1-\alpha}+z_{1-\alpha/2})}{\Delta}\Bigr)^2
  \;=\;\kappa(\alpha)\,n_{\mathrm{opt}},
\]
where throughout $z_{1-x}:=\Phi^{-1}(1-x)$ (see the notation note below).
Hence $n^\star_{3}\ge\kappa(\alpha)\,n_{\mathrm{opt}}$.

\emph{Step 2 (randomization).}  The budget-split argument of Step~1 is already stated for
arbitrary measurable commit profiles: replace the indicator of $A_\pm$ by a randomized test
function $\psi_\pm:\R\to[0,1]$, set $p:=\E_0\psi_+$, $q:=\E_0\psi_-$, and apply the randomized
Neyman--Pearson lemma on each side.  The extremal randomized test at a given $P_0$-level is again
a (possibly boundary-randomized) one-sided tail, so the same two inequalities hold verbatim and
the same minimization over $p+q\le\alpha$ applies.  No randomized rule beats the symmetric band.

\emph{Step 3 (impossibility at $4n_{\mathrm{opt}}$).}  At $n=4n_{\mathrm{opt}}$ one has
$\Delta/t=2z_{1-\alpha}$.  Step~1 requires $z_{1-\alpha}+z_{1-\alpha/2}\le2z_{1-\alpha}$, i.e.\
$z_{1-\alpha/2}\le z_{1-\alpha}$.  But $z_{1-\alpha/2}>z_{1-\alpha}$ for every $\alpha\in(0,\tfrac12)$,
so the inequality fails: no $\alpha$-certifying rule exists at $n=4n_{\mathrm{opt}}$.  Since
$\kappa(\alpha)=(1+z_{1-\alpha/2}/z_{1-\alpha})^2>4$ for every finite $\alpha$, the ratio
$\kappa(\alpha)/4>1$ cannot be closed to $1$ --- it is a theorem of impossibility, not an open
constant gap.
\end{proof}

\begin{corollary}[Status of the Wave~3 residual]\label{cor:t1c-status}
The tight-constants program is \textbf{closed} for the three-world Gaussian certifying problem:
the exact minimax constant is $\kappa(\alpha)$, not $4$ and not $16$.  The certificate's
two-world $4\times$ penalty (Remark~\ref{rmk:opt-delicate}) and the three-world floor
$4n_{\mathrm{opt}}$ (Prop.~\ref{prop:t1b-3world}) address \emph{distinct} constraints; only the
former is about the deployed certificate's CI width, while the latter is a loose pairwise bound
subsumed by Theorem~\ref{thm:t1c-exact}.
\end{corollary}

\begin{remark}[Numerical confirmation]\label{rmk:t1c-val}
\texttt{theory\_v2/val\_tight\_constants.py} (section \texttt{T1c}) verifies:
\textbf{[C1]} no two-threshold (symmetric or asymmetric) rule certifies at any
$n<(1-\varepsilon)\kappa(\alpha)n_{\mathrm{opt}}$ on a dense grid;
\textbf{[C2]} randomized mixture search over threshold rules does not beat $\kappa$;
\textbf{[C3]} at $n=4n_{\mathrm{opt}}$ the three-world LP is infeasible for every tested $\alpha$;
\textbf{[C4]} at $n=\kappa(\alpha)n_{\mathrm{opt}}$ the symmetric band binds all three errors at
$\alpha$.
\end{remark}
